\documentclass[10pt,a4paper]{amsart}
\usepackage{amsmath,amssymb,amsthm,graphicx}
% Retained preamble for the Stacks Project chapter {CH} (shared project preamble).
% The chapter file pinned-original-source.tex only carries \input{{preamble}} on its first line,
% so no preamble text lives inside the chapter itself.  The Stacks Project uses one single shared
% preamble for every chapter: the snapshot stages it three times (sources/stacks-chow/preamble.tex,
% sources/stacks-etale-cohomology/preamble.tex and focused-seeds/AGP/sources/stacks-snapshot-20261011/
% preamble.tex) and those copies are identical in size.  The block below is copied mechanically,
% line for line, from sources/stacks-chow/preamble.tex lines 14-18 and 154-217: the Xy-pic setup and
% the native theorem declarations together with the mathematical macros.  Nothing was added, removed
% or redefined; the wrapper class and the amsmath/amssymb/amsthm/graphicx base packages are supplied
% by the assembler.
% For commutative diagrams we use Xy-pic
% Chapter: stacks-spaces-more-morphisms.tex -- More on Morphisms of Spaces.
\usepackage[all]{xy}

% We use 2cell for 2-commutative diagrams.
\xyoption{2cell}

% Theorem environments.
%
\theoremstyle{plain}
\newtheorem{theorem}[subsection]{Theorem}
\newtheorem{proposition}[subsection]{Proposition}
\newtheorem{lemma}[subsection]{Lemma}

\theoremstyle{definition}
\newtheorem{definition}[subsection]{Definition}
\newtheorem{example}[subsection]{Example}
\newtheorem{exercise}[subsection]{Exercise}
\newtheorem{situation}[subsection]{Situation}

\theoremstyle{remark}
\newtheorem{remark}[subsection]{Remark}
\newtheorem{remarks}[subsection]{Remarks}

\numberwithin{equation}{subsection}

% Macros
%
\def\lim{\mathop{\mathrm{lim}}\nolimits}
\def\colim{\mathop{\mathrm{colim}}\nolimits}
\def\Spec{\mathop{\mathrm{Spec}}}
\def\Hom{\mathop{\mathrm{Hom}}\nolimits}
\def\Ext{\mathop{\mathrm{Ext}}\nolimits}
\def\SheafHom{\mathop{\mathcal{H}\!\mathit{om}}\nolimits}
\def\SheafExt{\mathop{\mathcal{E}\!\mathit{xt}}\nolimits}
\def\Sch{\mathit{Sch}}
\def\Mor{\mathop{\mathrm{Mor}}\nolimits}
\def\Ob{\mathop{\mathrm{Ob}}\nolimits}
\def\Sh{\mathop{\mathit{Sh}}\nolimits}
\def\NL{\mathop{N\!L}\nolimits}
\def\CH{\mathop{\mathrm{CH}}\nolimits}
\def\proetale{{pro\text{-}\acute{e}tale}}
\def\etale{{\acute{e}tale}}
\def\QCoh{\mathit{QCoh}}
\def\Ker{\mathop{\mathrm{Ker}}}
\def\Im{\mathop{\mathrm{Im}}}
\def\Coker{\mathop{\mathrm{Coker}}}
\def\Coim{\mathop{\mathrm{Coim}}}

% Boxtimes
%
\DeclareMathSymbol{\boxtimes}{\mathbin}{AMSa}{"02}

%
% Macros for moduli stacks/spaces
%
\def\QCohstack{\mathcal{QC}\!\mathit{oh}}
\def\Cohstack{\mathcal{C}\!\mathit{oh}}
\def\Spacesstack{\mathcal{S}\!\mathit{paces}}
\def\Quotfunctor{\mathrm{Quot}}
\def\Hilbfunctor{\mathrm{Hilb}}
\def\Curvesstack{\mathcal{C}\!\mathit{urves}}
\def\Polarizedstack{\mathcal{P}\!\mathit{olarized}}
\def\Complexesstack{\mathcal{C}\!\mathit{omplexes}}
% \Pic is the operator that assigns to X its picard group, usage \Pic(X)
% \Picardstack_{X/B} denotes the Picard stack of X over B
% \Picardfunctor_{X/B} denotes the Picard functor of X over B
\def\Pic{\mathop{\mathrm{Pic}}\nolimits}
\def\Picardstack{\mathcal{P}\!\mathit{ic}}
\def\Picardfunctor{\mathrm{Pic}}
\def\Deformationcategory{\mathcal{D}\!\mathit{ef}}
\newenvironment{reference}{}{}
\newenvironment{slogan}{}{}
\newenvironment{history}{}{}
\usepackage{cleveref}
\crefname{originallemma}{original lemma}{original lemmas}
\Crefname{originallemma}{Original lemma}{Original lemmas}
\makeatletter
\newlabel{more-algebra-lemma-base-change-comparison}{{0661}{source}{}{sourcepointer-1}{}}
\newlabel{more-algebra-lemma-base-change-comparison@cref}{{[originallemma][0][]0661}{[1][0][]source}{}{}{}}
\makeatother
\crefname{originallemma}{original lemma}{original lemmas}
\Crefname{originallemma}{Original lemma}{Original lemmas}
\makeatletter
\newlabel{more-algebra-lemma-flat-descent-tor-amplitude}{{068S}{source}{}{sourcepointer-2}{}}
\newlabel{more-algebra-lemma-flat-descent-tor-amplitude@cref}{{[originallemma][0][]068S}{[1][0][]source}{}{}{}}
\makeatother
\crefname{originallemma}{original lemma}{original lemmas}
\Crefname{originallemma}{Original lemma}{Original lemmas}
\makeatletter
\newlabel{sites-cohomology-lemma-tor-amplitude-stalk}{{0DJJ}{source}{}{sourcepointer-3}{}}
\newlabel{sites-cohomology-lemma-tor-amplitude-stalk@cref}{{[originallemma][0][]0DJJ}{[1][0][]source}{}{}{}}
\makeatother
\crefname{originallemma}{original lemma}{original lemmas}
\Crefname{originallemma}{Original lemma}{Original lemmas}
\makeatletter
\newlabel{spaces-perfect-lemma-tor-independence-and-tor-amplitude}{{0DKD}{source}{}{sourcepointer-4}{}}
\newlabel{spaces-perfect-lemma-tor-independence-and-tor-amplitude@cref}{{[originallemma][0][]0DKD}{[1][0][]source}{}{}{}}
\makeatother
\crefname{originallemma}{original lemma}{original lemmas}
\Crefname{originallemma}{Original lemma}{Original lemmas}
\makeatletter
\newlabel{spaces-properties-theorem-exactness-stalks}{{04K5}{source}{}{sourcepointer-5}{}}
\newlabel{spaces-properties-theorem-exactness-stalks@cref}{{[originallemma][0][]04K5}{[1][0][]source}{}{}{}}
\makeatother
\title{Torus amplitude of a quasi-coherent complex on a space descends along fpqc covers of the base}
\author{The Stacks Project Authors}
\date{stacks-project:spaces-more-morphisms.tex:local-snapshot-20261011; GNU FDL 1.2 or later}
\begin{document}
\maketitle
\noindent\textit{Editorial note.} Statement span is lines 13020-13035 of sources/stacks-spaces-more-morphisms/pinned-original-source.tex; proof span is lines 13037-13075, the complete author proof (37 body lines).\par\medskip
\noindent\textit{Editorial note.} Navigation stubs: 5 cited-but-not-included labels are supplied as non-load-bearing pointers with their pinned tags where the snapshot tag index has them.\par\medskip
\noindent\textit{Editorial note.} extra\_packages is empty on purpose: the retained preamble already loads xy with options and a second declaration clashes in the assembler wrapper.\par\medskip
\noindent\textit{Editorial note.} Text is retained verbatim, including the author's own typos and placeholders.\par\medskip
\noindent\textit{Editorial note.} One typographic slip in the displayed formula (O\_\{X\_, x\_i\}) is retained verbatim from the author source.\par\medskip
\noindent\textit{Allowed standard prerequisites (editorial).}\begin{itemize}
\item Explicit prerequisite boundary. Stalk formulas for tor amplitude on sites, the base-change comparison lemma for derived tensor products (More on Algebra, lemma-base-change-comparison), and flat descent of tor amplitude (lemma-flat-descent-tor-amplitude).
\end{itemize}\medskip
\ifdefined\hypertarget\hypertarget{sourcepointer-1}{}\fi
\noindent\textit{Source reference.} Original lemma 0661: sources/stacks-spaces-more-morphisms/pinned-original-source.tex, source lines 13075--13075. Cited by the retained statement or proof but not part of the collected spans; supplied only as a navigation pointer. The argument does not depend on its content.\par\medskip
\ifdefined\hypertarget\hypertarget{sourcepointer-2}{}\fi
\noindent\textit{Source reference.} Original lemma 068S: sources/stacks-spaces-more-morphisms/pinned-original-source.tex, source lines 13075--13075. Cited by the retained statement or proof but not part of the collected spans; supplied only as a navigation pointer. The argument does not depend on its content.\par\medskip
\ifdefined\hypertarget\hypertarget{sourcepointer-3}{}\fi
\noindent\textit{Source reference.} Original lemma 0DJJ: sources/stacks-spaces-more-morphisms/pinned-original-source.tex, source lines 13075--13075. Cited by the retained statement or proof but not part of the collected spans; supplied only as a navigation pointer. The argument does not depend on its content.\par\medskip
\ifdefined\hypertarget\hypertarget{sourcepointer-4}{}\fi
\noindent\textit{Source reference.} Original lemma 0DKD: sources/stacks-spaces-more-morphisms/pinned-original-source.tex, source lines 13075--13075. Cited by the retained statement or proof but not part of the collected spans; supplied only as a navigation pointer. The argument does not depend on its content.\par\medskip
\ifdefined\hypertarget\hypertarget{sourcepointer-5}{}\fi
\noindent\textit{Source reference.} Original lemma 04K5: sources/stacks-spaces-more-morphisms/pinned-original-source.tex, source lines 13075--13075. Cited by the retained statement or proof but not part of the collected spans; supplied only as a navigation pointer. The argument does not depend on its content.\par\medskip
\begin{lemma}
\label{lemma-tor-amplitude-descends-fppf}
Let $S$ be a scheme. Let $\{g_i : Y_i \to Y\}$ be an fpqc covering of
algebraic spaces over $S$. Let $f : X \to Y$ be a morphism of algebraic
spaces and set $X_i = Y_i \times_Y X$ with projections $f_i : X_i \to Y_i$
and $g'_i : X_i \to X$. Let $E \in D_\QCoh(\mathcal{O}_X)$.
Let $a, b \in \mathbf{Z}$. Then the following are equivalent
\begin{enumerate}
\item $E$ has tor amplitude in $[a, b]$ as an object of
$D(f^{-1}\mathcal{O}_Y)$, and
\item $L(g'_i)^*E$ has tor amplitude in $[a, b]$ as a object of
$D(f_i^{-1}\mathcal{O}_{Y_i})$ for all $i$.
\end{enumerate}
Also true if ``tor amplitude in $[a, b]$'' is replaced by
``locally finite tor dimension''.
\end{lemma}

\begin{proof}
Pullback preserves ``tor amplitude in $[a, b]$'' by
Derived Categories of Spaces, Lemma
\ref{spaces-perfect-lemma-tor-independence-and-tor-amplitude}
Observe that $Y_i$ and $X$ are tor independent over $Y$
as $Y_i \to Y$ is flat. Let us assume (2) and prove (1).
We can compute tor dimension at stalks, see
Cohomology on Sites, Lemma \ref{sites-cohomology-lemma-tor-amplitude-stalk}
and Properties of Spaces, Theorem
\ref{spaces-properties-theorem-exactness-stalks}.
Let $\overline{x}$ be a geometric point of $X$. Choose
an $i$ and a geometric point $\overline{x}_i$ in $X_i$ with image
$\overline{x}$ in $X$. Then
$$
(L(g_i')^*E)_{\overline{x}_i} =
E_{\overline{x}}
\otimes_{\mathcal{O}_{X, \overline{x}}}^\mathbf{L}
\mathcal{O}_{X_, \overline{x}_i}
$$
Let $\overline{y}_i$ in $Y_i$ and $\overline{y}$ in $Y$
be the image of $\overline{x}_i$ and $\overline{x}$.
Since $X$ and $Y_i$ are tor independent over $Y$, we can apply
More on Algebra, Lemma \ref{more-algebra-lemma-base-change-comparison}
to see that the right hand side of the displayed formula is equal to
$E_{\overline{x}}
\otimes_{\mathcal{O}_{Y, \overline{y}}}^\mathbf{L}
\mathcal{O}_{Y_i, \overline{y}_i}$
in $D(\mathcal{O}_{Y_i, \overline{y}_i})$.
Since we have assume the tor amplitude of this is in
$[a, b]$, we conclude that the tor amplitude of
$E_{\overline{x}}$ in $D(\mathcal{O}_{Y, \overline{y}})$
is in $[a, b]$ by More on Algebra, Lemma
\ref{more-algebra-lemma-flat-descent-tor-amplitude}.
Thus (1) follows.

\medskip\noindent
Using some elementary topology the case
``locally finite tor dimension'' follows too.
\end{proof}

\end{document}
